\documentclass{report}
\usepackage{amsmath,amssymb}
\setlength{\parindent}{0mm}
\setlength{\parskip}{1em}
\begin{document}
\begin{center}
\rule{6in}{1pt} \
{\large Robert Vaselaar \\
{\bf Trigonometric Space-Time Discretization of the Gauge-Free Dirac Equation and its Implementation}}

Department of Mathematics and Statistics \\ South Dakota State University \\ Brookings \\ SD 57007
\\
{\tt robert.vaselaar@sdstate.edu}\\
Hyun Lim\\
Jung-Han Kimn\\
Dongming Mei\end{center}

We use a generalized space-time finite element method to find a new
computational approach to the time-based Dirac equation. In our method we
replace the standard $C^0$ linear Lagrangian elements with C1
trigonometric elements. $C^1$ elements preserve the continuity of
position, momentum and energy in the quantum wave and trigonometric basis
functions are conceptually closer to the function space of available
analytic solutions. Instead of using a space-only discretization with a
time-stepping scheme, we simultaneously discretize both the spatial and
temporal dimensions.

The result is a discrete form of the Dirac equation that converges to the
analytic solution without modifying the original Dirac operator.
Numerical results using PETSc and its KSP solvers in Cluster machines
will be presented with a possible parallelization implementation. Finally
we show that because we use a space-time discretization technique, our
method demonstrates the same Lorentz covariance as the continuum Dirac
equation.


\end{document}
